\documentclass[UTF8,a4paper,11pt,fontset=fandol]{ctexart}
\usepackage[margin=2.2cm]{geometry}
\usepackage{amsmath,amssymb,booktabs}
\setlength{\parskip}{0.3em}
\linespread{1.1}

\begin{document}
\pagestyle{plain}

\begin{center}
  {\Large\bfseries 第二次上机作业：四种 QR 分解方法的正交性比较}
\end{center}

\section{原题抄录}

随机生成 10 个不同条件数的 $500\times20$ 阶矩阵，编程比较
CGS、MGS、MGS2 和 Householder 变换进行 QR 分解后得到的
$\|Q^{\mathrm T}Q-I\|_{\mathrm F}$，以衡量计算所得 $Q$ 的正交性。
要求程序输出这 10 个矩阵分别使用四种方法得到的范数。

本文按题目给定公式计算。严格地说，
$\|Q^{\mathrm T}Q-I\|_{\mathrm F}$ 是正交性误差，未进行相对归一化；
若定义相对误差，则还需除以 $\|I\|_{\mathrm F}=\sqrt{20}$。
本实验不增加这一除法，以保持与所给程序一致。

\section{代码的结构、算法说明}

\subsection*{（1）矩阵与误差的记号}

记待分解矩阵为
\[
A=[a_1,\ldots,a_n]\in\mathbb R^{m\times n},
\qquad m=500,\quad n=20.
\]
对列满秩矩阵，薄 QR 分解写成
\[
A=QR,\qquad
Q=[q_1,\ldots,q_n]\in\mathbb R^{m\times n},
\qquad R=(r_{ij})\in\mathbb R^{n\times n},
\]
其中 $R$ 为上三角矩阵，精确算术下 $Q^{\mathrm T}Q=I_n$。
程序使用双精度浮点数计算，所得 $Q$ 可能失去正交性。

函数 \texttt{orthogonality\_error(Q)} 计算
\[
\eta(Q)=\|Q^{\mathrm T}Q-I_n\|_{\mathrm F}
=\left(
\sum_{i=1}^{n}\sum_{j=1}^{n}
\bigl(q_i^{\mathrm T}q_j-\delta_{ij}\bigr)^2
\right)^{1/2},
\]
其中 $\delta_{ij}$ 为 Kronecker 符号。
$\eta(Q)$ 越小，表示所得列向量越接近标准正交向量组。
它没有必须小于 1 的要求。

\subsection*{（2）随机矩阵的生成}

函数
\texttt{generate\_matrix(m, n, condition\_number, rng)}
先生成两个标准正态随机矩阵，再调用
\texttt{numpy.linalg.qr}
得到列正交矩阵 $U\in\mathbb R^{500\times20}$ 和正交矩阵
$V\in\mathbb R^{20\times20}$。

设目标条件数为 $\kappa$，取
\[
\sigma_i=\kappa^{-(i-1)/(n-1)},
\qquad i=1,\ldots,n,
\]
并令
\[
\Sigma=\operatorname{diag}(\sigma_1,\ldots,\sigma_n),
\qquad A=U\Sigma V^{\mathrm T}.
\]
精确算术下，$A$ 的奇异值为
$\sigma_1,\ldots,\sigma_n$，因此
\[
\kappa_2(A)=\frac{\sigma_1}{\sigma_n}=\kappa.
\]

主程序使用 \texttt{numpy.logspace(1, 16, 10)}
选取 10 个目标条件数，即
\[
\kappa_j=10^{1+15(j-1)/9},
\qquad j=1,\ldots,10.
\]
固定随机种子为 2026，每次生成矩阵时使用新的随机数。
程序还通过
\texttt{numpy.linalg.}\allowbreak\texttt{cond(A)}
计算实际二范数条件数，用于观察矩阵构造中的浮点舍入影响。

这里调用库函数 QR 仅用于生成测试矩阵；
下面参与比较的四种 QR 算法均由程序直接实现。
对每个测试矩阵，四种算法使用相同输入，
且各函数均不修改原始矩阵 $A$。

\subsection*{（3）经典 Gram--Schmidt 方法：CGS}

函数 \texttt{CGS(A)} 逐列计算 $Q,R$。
处理第 $j$ 列时，先计算
\[
r_{ij}=q_i^{\mathrm T}a_j,
\qquad i=1,\ldots,j-1,
\]
再减去全部投影：
\[
v=a_j-\sum_{i=1}^{j-1}r_{ij}q_i.
\]
最后令
\[
r_{jj}=\|v\|_2,
\qquad q_j=\frac{v}{r_{jj}}.
\]

当 $j=1$ 时求和为空。
代码先用一个循环计算全部投影系数，
再用另一个循环减去投影；
各系数始终由原始列 $a_j$ 计算。

\subsection*{（4）改进 Gram--Schmidt 方法：MGS}

函数 \texttt{MGS(A)} 处理第 $j$ 列时，
先令 $v=a_j$，然后对 $i=1,\ldots,j-1$ 依次执行
\[
r_{ij}=q_i^{\mathrm T}v,
\qquad
v\leftarrow v-r_{ij}q_i.
\]
完成投影消去后令
\[
r_{jj}=\|v\|_2,
\qquad q_j=\frac{v}{r_{jj}}.
\]

与 CGS 不同，每次计算内积时使用的是当前残差 $v$。
两种方法在精确算术下等价，但浮点运算顺序不同，
因而数值稳定性不同。

\subsection*{（5）两遍重正交化方法：MGS2}

函数 \texttt{MGS2(A)} 对每一列连续执行两遍
MGS 投影消去，最后再归一化。

处理第 $j$ 列时，先令
\[
v=a_j,\qquad r_{ij}=0\quad(i<j).
\]
对于遍数 $p=1,2$，
每遍按 $i=1,\ldots,j-1$ 依次执行
\[
c_{ij}^{(p)}=q_i^{\mathrm T}v,
\]
\[
v\leftarrow v-c_{ij}^{(p)}q_i,
\qquad
r_{ij}\leftarrow r_{ij}+c_{ij}^{(p)}.
\]
两遍结束后，令
\[
r_{jj}=\|v\|_2,
\qquad q_j=\frac{v}{r_{jj}}.
\]

因此
\[
r_{ij}=c_{ij}^{(1)}+c_{ij}^{(2)}.
\]
第二遍用于消除第一遍因舍入误差而残留的投影；
其系数需要累加到 $R$ 中，不能覆盖第一遍的结果。

\subsection*{（6）Householder QR 方法}

函数 \texttt{Householder(A)} 采用教材中的记号
\[
H=I-2\frac{vv^{\mathrm T}}{v^{\mathrm T}v},
\qquad v\ne0.
\]
若 $\|x\|_2=\|y\|_2$ 且 $x\ne y$，
取 $v=x-y$，则 $Hx=y$。

令 $A_0=A$。第 $k$ 步取
\[
x=A_{k-1}(k:m,k)\in\mathbb R^{m-k+1},
\qquad k=1,\ldots,n.
\]
定义
\[
\operatorname{sgn}_+(t)=
\begin{cases}
1,&t\ge0,\\
-1,&t<0,
\end{cases}
\]
并取
\[
y=-\operatorname{sgn}_+(x_1)\|x\|_2e_1,
\qquad v=x-y,
\]
其中 $e_1\in\mathbb R^{m-k+1}$ 为第一个标准基向量。
这一符号选择使 $v$ 的第一个分量由同号相加得到，
避免相减消去。

将局部反射向量 $v$ 补零为
\[
v_k=
\begin{bmatrix}
0_{k-1}\\v
\end{bmatrix},
\qquad
H_k=I_m-2\frac{v_kv_k^{\mathrm T}}{v_k^{\mathrm T}v_k},
\]
并计算
\[
A_k=H_kA_{k-1},
\]
从而消去第 $k$ 列对角线以下的元素。

程序不显式构造 $H_k$，
只对当前右下角子矩阵 $B$ 执行
\[
B\leftarrow B-2\frac{v(v^{\mathrm T}B)}{v^{\mathrm T}v}.
\]
代码中的 \texttt{np.outer(v, v @ B)}
对应 $v(v^{\mathrm T}B)$。
每一步的局部向量 $v$
保存到列表 \texttt{vectors} 中。

完成全部反射后，理论上有
\[
A_n=H_n\cdots H_1A
=
\begin{bmatrix}
R\\0_{(m-n)\times n}
\end{bmatrix}.
\]
代码取前 $n$ 行，再用 \texttt{np.triu}
清除下三角的舍入残留，得到 $R$。

由 $H_k^{\mathrm T}=H_k^{-1}=H_k$，薄正交因子为
\[
Q=H_1H_2\cdots H_nE,
\qquad
E=
\begin{bmatrix}
I_n\\0_{(m-n)\times n}
\end{bmatrix}.
\]
代码先用 \texttt{np.eye(m, n)} 初始化 $E$，
再按 $n,n-1,\ldots,1$ 的顺序逐次左乘反射矩阵，
得到上述乘积。

数学下标从 1 开始，而 Python 下标从 0 开始。

\subsection*{（7）主程序与输出}

主程序对每个矩阵依次调用
\texttt{CGS}、\texttt{MGS}、
\texttt{MGS2} 和 \texttt{Householder}，
各函数返回薄矩阵 $Q$ 和上三角矩阵 $R$。
随后计算四个正交性误差，并以科学计数法输出。

本次比较的指标仅为
$\|Q^{\mathrm T}Q-I_n\|_{\mathrm F}$，
不包含运行时间和分解残差
$\|A-QR\|_{\mathrm F}/\|A\|_{\mathrm F}$。

\section{输出结果呈现}

以下为所给程序的一次实际运行结果，
使用 Python 3.12、NumPy 2.3.5，
并将 OpenBLAS 线程数设为 1。
不同数值库、线程设置或平台可能改变舍入误差，
因此固定随机种子并不保证跨平台输出逐位相同。

为便于排版，将程序输出拆为两个表，
以相同序号对应。条件数保留程序输出的小数位数。

\begin{center}
\small
\begin{tabular}{crr}
\toprule
序号 & 目标条件数 & 实际条件数 $\kappa_2(A)$\\
\midrule
1  & \texttt{1.0000e+01} & \texttt{1.0000e+01}\\
2  & \texttt{4.6416e+02} & \texttt{4.6416e+02}\\
3  & \texttt{2.1544e+04} & \texttt{2.1544e+04}\\
4  & \texttt{1.0000e+06} & \texttt{1.0000e+06}\\
5  & \texttt{4.6416e+07} & \texttt{4.6416e+07}\\
6  & \texttt{2.1544e+09} & \texttt{2.1544e+09}\\
7  & \texttt{1.0000e+11} & \texttt{1.0000e+11}\\
8  & \texttt{4.6416e+12} & \texttt{4.6416e+12}\\
9  & \texttt{2.1544e+14} & \texttt{2.1519e+14}\\
10 & \texttt{1.0000e+16} & \texttt{9.0691e+15}\\
\bottomrule
\end{tabular}
\end{center}

四种方法的正交性误差 $\eta(Q)$ 如下：

\begin{center}
\small
\setlength{\tabcolsep}{5pt}
\begin{tabular}{crrrr}
\toprule
序号 & CGS & MGS & MGS2 & Householder\\
\midrule
1
& \texttt{3.478256e-15}
& \texttt{2.705890e-15}
& \texttt{1.065330e-15}
& \texttt{2.310313e-15}\\
2
& \texttt{1.688335e-12}
& \texttt{6.547062e-14}
& \texttt{1.643985e-15}
& \texttt{2.036919e-15}\\
3
& \texttt{6.940720e-09}
& \texttt{1.581109e-12}
& \texttt{1.141623e-15}
& \texttt{1.947947e-15}\\
4
& \texttt{4.890503e-05}
& \texttt{3.564626e-11}
& \texttt{1.601224e-15}
& \texttt{2.452404e-15}\\
5
& \texttt{2.190315e-02}
& \texttt{3.412914e-09}
& \texttt{1.487992e-15}
& \texttt{2.168266e-15}\\
6
& \texttt{2.707436e+00}
& \texttt{1.324775e-07}
& \texttt{1.732023e-15}
& \texttt{1.774602e-15}\\
7
& \texttt{6.767364e+00}
& \texttt{4.111114e-06}
& \texttt{1.465112e-15}
& \texttt{2.293422e-15}\\
8
& \texttt{7.802343e+00}
& \texttt{2.169458e-04}
& \texttt{1.129266e-15}
& \texttt{2.232910e-15}\\
9
& \texttt{8.400029e+00}
& \texttt{6.831170e-03}
& \texttt{1.573747e-15}
& \texttt{2.303039e-15}\\
10
& \texttt{1.024099e+01}
& \texttt{3.842240e-01}
& \texttt{1.393481e-15}
& \texttt{1.969044e-15}\\
\bottomrule
\end{tabular}
\end{center}

\section{对输出结果的分析和结论}

\subsection*{（1）CGS 与 MGS 的比较}

当条件数较小时，四种方法均能得到接近标准正交的列向量。
随着本次测试矩阵条件数增大，
CGS 的正交性误差增长明显。
例如，条件数约为 $10^6$ 时，
其误差约为 $4.89\times10^{-5}$；
条件数约为 $2.15\times10^9$ 时，
误差已超过 1，说明正交性严重丢失。

MGS 的结果明显优于 CGS。
例如，在上述两个矩阵上，MGS 的误差分别约为
$3.56\times10^{-11}$ 和 $1.32\times10^{-7}$。
这是因为 MGS 每次对更新后的残差进行投影，
减轻了舍入误差的影响。

但 MGS 也不能完全避免病态矩阵带来的正交性损失：
最后一个矩阵的 MGS 误差约为 $3.84\times10^{-1}$。

\subsection*{（2）MGS2 与 Householder 的表现}

本次 10 个矩阵上，MGS2 的误差均约为 $10^{-15}$。
第二遍投影消去有效去除了第一遍留下的非正交分量。
Householder 的误差也始终约为 $10^{-15}$，
表现出较好的正交性保持能力。
这与它通过一系列正交反射构造 $Q$ 的特点一致。

双精度浮点数的机器精度约为 $2.22\times10^{-16}$。
本实验计算的是包含多个分量的 Frobenius 范数，
误差还受到矩阵维数和运算累积影响，
因此 $10^{-15}$ 量级与接近机器精度的正交性相符。

本次数据中 MGS2 的误差略小于 Householder，
但两者处于相同数量级，
不能据此认定 MGS2 在所有矩阵上都优于 Householder。

\subsection*{（3）实际条件数与实验范围}

前几组矩阵的实际条件数与目标值在显示精度内一致。
当目标条件数达到 $10^{16}$ 时，
实际计算值约为 $9.07\times10^{15}$。
此时预设最小奇异值为 $10^{-16}$，
已接近双精度舍入误差尺度，
构造矩阵和计算奇异值时的舍入都会影响实际条件数。
因此，高条件数组应结合实际条件数解释，
不能将目标值视为精确实现值。

本实验每个目标条件数只取一个随机矩阵，
结论反映这组样本的表现，
并非对全部同条件数矩阵的统计结论。

此外，正交性误差小不单独证明 $A\approx QR$；
若要全面评价分解质量，
还应另行检查相对分解残差。
本次没有测量运行时间，
因此不对四种方法的实际耗时作排序。

\subsection*{（4）结论}

对本次生成的 10 个测试矩阵，
随着条件数增大，CGS 较早出现明显的正交性损失；
MGS 有所改善，但在极高条件数下仍失去正交性。
MGS2 和 Householder 在全部样本上均保持约
$10^{-15}$ 量级的正交性误差。

因此，在本次实验范围内，
重正交化与 Householder 反射都能有效改善
计算所得 $Q$ 的正交性。

\end{document}